Modern neuroscience lacks robust methods to characterise long-duration and
continual time series, especially in settings where the statistical properties
of the data evolve over time. This limitation present a methodological gap that
must be addressed in order to unlock the scientific potential of NaLoDuCo
experiments.

We will address these challenges by developing a new NaLoDuCo-focused
environment of machine-learning analysis methods specifically tailored to
long-duration and continual time series.

Our first goal will be to provide implementations of methods already used at
the GCNU, SWC, and AIND adapted to optimally process long-duration time series.
These methods span multiple domains—kinematics, neural dynamics, behavioural
state segmentation, forecasting, and joint modelling—and are designed to work
together within an integrated analysis pipeline.

Below we first summarise the functionality of this pipeline, then we highlight
challenges encountered when using it in long-duration behavioral and neural
time series, and finally describe methods that we will use to address these
challenges.

\subsubsubsection*{Pipeline of methods used at the GCNU, SWC, and AIND to
process behavioral and neural time series}

\noindent\textbf{Behavioural Analysis:}

The first step in behavioural analysis involves multi-body-part tracking. For
this, we will use deep learning-based pose estimation methods such as
\href{https://github.com/talmolab/sleap}{SLEAP}, which enable accurate and
efficient tracking of multiple unmarked mice across long recording sessions.

From the tracked poses, we will infer continuous kinematic variables using
linear dynamical systems (LDS), including particle filter-based variants to
handle uncertainty and measurement noise. These kinematic features will be used
to infer discrete behavioural states with Hidden Markov Models (HMMs), as
implemented in tools such as
\href{https://dattalab.github.io/moseq2-website/index.html}{MoSeq}.

We will relate these inferred states and kinematic variables to
foraging-related outcomes—such as patch-leaving decisions—using both
generalised linear models (GLMs) and deep neural networks. These models will
allow us to capture both interpretable and high-capacity representations of
behavioural decision-making processes.

To recover the latent strategies guiding animal behaviour, we will apply
inverse reinforcement learning methods such as
\href{https://github.com/haozhu10015/hiql}{HIQL}, which estimate the underlying
reward functions and policies based on observed actions.

NaLoDuCo recordings uniquely support behavioural forecasting over extended
horizons—ranging from hours to days—far beyond what is feasible in conventional
short-duration experiments. To capitalise on this, we will apply long-horizon
forecasting models using recurrent neural networks (RNNs) and transformer
architectures, which are well-suited to modelling long-range temporal
dependencies.

\vspace{1em}
\noindent\textbf{Neural Data Analysis:}

Analysis of high-density electrophysiology will begin with latent variable
modelling to reduce the dimensionality of population neural recordings. We will
use both linear and nonlinear latent dynamics models, including
\href{https://github.com/joacorapela/svGPFA}{svGPFA}, which uses Gaussian
processes, and \href{https://snel.ai/resources/lfads/}{LFADS}, a deep
generative model based on recurrent neural networks.

The resulting low-dimensional trajectories will be used to infer discrete
neural states via HMMs, using methods such as
\href{https://github.com/lindermanlab/ssm}{SSM}. For neural activity
forecasting across long duration, we will again employ RNNs and transformers,
which can model complex temporal structure in spiking activity.

We will also decode the animal’s position from hippocampal spike trains using
point-process decoders, enabling the analysis of spatial coding and replay
phenomena during naturalistic foraging behaviour. We will build on existing
implementations such as
\href{https://github.com/Eden-Kramer-Lab/replay_trajectory_classification}{replay\_trajectory\_classification}.

\vspace{1em}
\noindent\textbf{Joint Neural-Behavioural Modelling:}  

To understand the interactions between neural dynamics and behaviour, we will
use models that extract shared latent representations from both domains. These
models will help reveal how cognitive and behavioural states are jointly
encoded in neural activity.

We will adapt Recognition-Parameterised Models (RPM), a Bayesian approach
developed at the GCNU, which infers latent variables that explain multiple
observation streams through highly nonlinear relationships. We will also use
\href{https://cebra.ai/}{CEBRA}, a state-of-the-art contrastive learning
framework designed for multimodal representation learning, to discover
temporally and semantically aligned neural-behavioural structure.

\begin{table}
    \caption{Initial data analysis methods to disseminate}
    \label{table:initialMethodsToDisseminate}
    \begin{longtable}{|p{1.9cm}|p{3cm}|p{2.5cm}|p{4.5cm}|}
        \hline
        \textbf{Domain} & \textbf{Functionality} & \textbf{Method} & \textbf{Model Type} \\
        \hline\hline
        behaviour & multi-body-part tracking & SLEAP & deep neural network \\
        \hline
        behaviour & kinematics inference & LDS & linear dynamical system \\
        \hline
        behaviour & kinematics inference & LDS & particle filter \\
        \hline
        behaviour & state inference & SSM & hidden Markov model \\
        \hline
        behaviour & regression &  & generalised linear model \\
        \hline
        behaviour & regression &  & deep neural network \\
        \hline
        behaviour & policy inference & L(M)V-IQL & reinforcement learning \\
        \hline
        behaviour & long-duration forecasting &  & RNN \\
        \hline
        behaviour & long-duration forecasting &  & transformers \\
        \hline\hline
        brain & latents inference & svGPFA & Gaussian processes \\
        \hline
        brain & latents inference & LFADS & RNN \\
        \hline
        brain & state inference & SSM & hidden Markov model \\
        \hline
        brain & long-duration forecasting &  & RNN \\
        \hline
        brain & long-duration forecasting &  & transformers \\
        \hline
        brain & decoding & NA & point-process decoder \\
        \hline\hline
        brain \& behaviour & latents inference & RPM & Bayesian inference + deep neural network \\
        \hline
        brain \& behaviour & latents inference & CEBRA & contrastive learning \\
        \hline
    \end{longtable}
\end{table}

\subsubsubsection*{Challenges using the previous pipeline in NaLoDuCo
experiments}

\subsubsubsection*{Non-stationarity}

Many methods for analysing neural and behavioural time series assume that the
underlying data-generating processes are stationary—that is, their statistical
properties remain constant over time. While this assumption may be acceptable
in short-duration experiments, it breaks down in long-duration and continual
recordings. In such settings, animals learn and adapt, their internal states
and motivations fluctuate, and their behaviour and physiology are influenced by
biological rhythms such as circadian, ultradian, and infradian cycles. These
changes induce non-stationarity in the data, making models that assume
stationarity progressively less reliable or even obsolete.

To address this challenge, we will adapt and develop methods that are
explicitly designed to operate in non-stationary environments. Our approach
draws on techniques from multiple domains, including adaptive signal
processing, machine learning, and Bayesian inference.

\paragraph{Adaptive Signal Processing.}

The field of adaptive signal processing has produced robust methods for
modelling linear systems with time-varying dynamics~\cite{haykin02}. The
recursive least-squares (RLS) algorithm, for example, is a well-known
adaptation of the ordinary least squares algorithm that continuously updates
model parameters to perform linear regression under non-stationary conditions.
We will use RLS to study time-varying relations between behavioural states, as
inferred by hidden Markov models, and foraging visit duration.

\paragraph{Adaptive State-Space Models.}

In state-space modelling, the Kalman filter provides a principled way to handle
non-stationary Gaussian processes with drifting mean and covariance. More
flexible approaches are needed when data exhibit abrupt regime shifts or
complex latent dynamics. Switching state-space models, such as Switching Linear
Dynamical Systems (SLDS) and Switching Hidden Markov Models (sHMMs), model
discrete changes in underlying system dynamics. For nonlinear, non-Gaussian
signals, particle filters approximate the posterior distribution through
sequential sampling. Bayesian online learning techniques offer a general
framework for continually updating model parameters as new data arrive. Using
these techniques we will build models that robustly infer kinematics over
months.

\paragraph{Non-stationarity in machine learning.}

Conventional machine-learning methods train models using schemes such as
stochastic gradient descent (SGD), which implicitly assume stationarity in the
data.  Online variants (similar to RLS, above) allow for some adaptation, but
in deep non-linear systems simple online updates may lead to loss of plasticity
preventing the trained model from tracking changes in the
data~\cite{ashAndAdams20}.

Instead, machine-learning approaches rely on a combination of (1) detection
methods, which monitor for significant changes in data distribution; (2)
adaptation methods, which incrementally update models using strategies such as
sliding windows, online learning, or ensemble-based approaches; and (3)
forgetting mechanisms, which allow models to discard outdated information while
retaining relevant past knowledge.  Our own recent work has shown how these
schemes may be blended into a single coherent learning
approach~\cite{galashovEtAl24}.

We will adapt these and similar approaches to extend many of the non-linear
base methods of Table 1 to handle drift in data statistics. In particular, we
have had previous success with non-stationarity in image classification and
online reinforcement learning settings~\cite{galashovEtAl24}.  We will extend
these to models that capture more complex relationships than simple
input-output mappings, including transformers,
\href{https://github.com/gatsby-sahani/rpm-aistats-2023}{Recognition-Parameterised
Models (RPMs)}, contrastive models including CEBRA and others listed above.

In summary, robust analysis of NaLoDuCo datasets requires models that
continuously adapt to evolving data distributions. Our offline analysis
framework will integrate both established adaptive algorithms and cutting-edge
methods from continual learning and concept drift to meet this challenge.

\subsubsubsection*{Computational efficiency and scaling}

The large dataset sizes generated by NaLoDuCo experimentation pose a
significant challenge for fast data analysis. We will develop accelerated
implementations of all methods in the library.  These implementations will use
\href{https://docs.jax.dev/}{JAX} for model learning, inference, and numerical
computation, \href{https://spark.apache.org/}{Apache Spark} or
\href{https://www.dask.org/}{Dask} to distribute pre-processing and feature
extraction, and Ray to distribute machine learning and deep learning
functionality.

\subsubsubsection{Deliverables}

\begin{enumerate}

	\item repository containing implementations of machine learning algorithms
	for offline processing NaLoDuCo experimental data, adapted to operate in
	non-stationary environments, and optimised to perform at scale when running
	on public clouds or institutional high-performance-computing systems.

	\item deployment of the methods in 1 in Amazon EC2 instances, to allow
	users to analyse on the cloud data from previous SWC NaLoDuCo foraging
	datasets and AIND NaLoDuCo odour learning datasets stored in the DANDI
	respository.

\end{enumerate}
